\section{Paragraf 16-Byte, Struktur Segmen, dan Overlapping}
\label{sec:paragraf-overlap}

Perkalian dengan \hex{10} pada rumus alamat fisik bukan pilihan
sewenang-wenang. Ia adalah akibat langsung dari cara memori
dibagi-bagi. Bagian ini menjelaskan satuan pembagian tersebut, lalu
menunjukkan akibatnya: satu alamat fisik dapat dinyatakan oleh
banyak pasangan \eng{segment:offset}.

\subsection{Paragraf 16-Byte}

Memori 8086 dibagi menjadi satuan berukuran 16 byte yang disebut
\term{paragraf}. Setiap paragraf diberi nomor urut mulai dari nol.
Paragraf pertama berisi byte \hex{00000} sampai \hex{0000F},
paragraf kedua berisi byte \hex{00010} sampai \hex{0001F}, dan
seterusnya. Satu paragraf selalu dimulai pada alamat yang digit
heksadesimal terakhirnya nol.

\begin{definisi}
\textbf{Paragraf} (\eng{paragraph}) adalah blok memori berukuran
16 byte yang menjadi satuan dasar penomoran segmen. Nilai sebuah
segmen adalah nomor paragraf tempat segmen itu dimulai. Karena itu
alamat dasar sebuah segmen selalu merupakan kelipatan 16 dan selalu
berakhir dengan digit \hex{0}.
\end{definisi}

Nomor paragraf dinyatakan dengan 16 bit, sehingga jumlah paragraf
yang dapat ditunjuk adalah \dec{65.536} buah. Karena satu paragraf
berisi 16 byte, seluruh paragraf tersebut meliputi
$65536 \times 16 = \dec{1.048.576}$ byte, yaitu tepat 1 MB. Dengan
demikian, himpunan alamat dasar segmen yang mungkin tepat menutupi
seluruh ruang alamat 8086. Inilah sebabnya nilai segmen yang sah
selalu berada pada rentang \hex{0000} sampai \hex{FFFF}.

\subsection{Struktur Segmen 64 KB}

Sebuah segmen bukanlah blok memori yang secara fisik terpisah,
melainkan sebuah \emph{jendela} 64 KB yang dimulai pada salah satu
alamat dasar paragraf. Offset 16-bit di dalam segmen bergerak dari
\hex{0000} sampai \hex{FFFF}. Bila dasar segmen adalah $B$, maka
segmen tersebut mencakup alamat $B$ sampai $B + \hex{FFFF}$, yaitu
\dec{65.536} byte berturut-turut.

Perhatikan bahwa panjang jendela tersebut hanya 64 KB, sedangkan
jarak antaralamat dasar dua segmen yang berurutan hanyalah 16 byte.
Akibatnya, jendela segmen \hex{0000} dan jendela segmen \hex{0001}
saling menumpuk hampir seluruhnya: hanya 16 byte pertama jendela
pertama yang tidak tercakup jendela kedua. Penumpukan inilah yang
disebut \term{overlapping segment}, dan sifatnya melekat pada
rancangan 8086, bukan sebuah kesalahan.

\subsection{Satu Alamat Fisik, Banyak Nama}

Karena setiap segmen hanya berjarak 16 byte dari segmen berikutnya,
sebuah alamat fisik dapat dibentuk oleh banyak pasangan yang
berbeda. Sebagai contoh, alamat fisik \hex{00050} dapat dinyatakan
oleh enam pasangan seperti tampak pada
Tabel~\ref{tab:overlap-pasangan}.

\begin{table}[H]
    \centering
    \caption{Enam pasangan \eng{segment:offset} yang menunjuk alamat
    fisik \hex{00050}.}
    \label{tab:overlap-pasangan}
    \small
    \begin{tabular}{@{}cp{1.6cm}p{1.6cm}p{3.9cm}p{2.1cm}@{}}
        \toprule
        \textbf{No} & \textbf{Segmen} & \textbf{Offset} &
        \textbf{Perhitungan} & \textbf{Alamat fisik} \\
        \midrule
        1 & \hex{0000} & \hex{0050} &
        \hex{00000} + \hex{0050} & \hex{00050} \\
        \addlinespace
        2 & \hex{0001} & \hex{0040} &
        \hex{00010} + \hex{0040} & \hex{00050} \\
        \addlinespace
        3 & \hex{0002} & \hex{0030} &
        \hex{00020} + \hex{0030} & \hex{00050} \\
        \addlinespace
        4 & \hex{0003} & \hex{0020} &
        \hex{00030} + \hex{0020} & \hex{00050} \\
        \addlinespace
        5 & \hex{0004} & \hex{0010} &
        \hex{00040} + \hex{0010} & \hex{00050} \\
        \addlinespace
        6 & \hex{0005} & \hex{0000} &
        \hex{00050} + \hex{0000} & \hex{00050} \\
        \bottomrule
    \end{tabular}
\end{table}

Pola pada tabel tersebut mudah dibaca: setiap kali nomor segmen
bertambah satu, dasar segmen bertambah 16, sehingga offset harus
berkurang 16 agar jumlahnya tetap. Untuk alamat fisik $F$ yang
berada di bawah \hex{10000}, banyaknya pasangan yang mungkin sama
dengan $\lfloor F / 16 \rfloor + 1$. Pada contoh di atas,
$\lfloor \hex{50} / \hex{10} \rfloor + 1 = 5 + 1 = 6$ pasangan.

Gambar~\ref{fig:overlap-segmen} memperlihatkan hal yang sama dalam
bentuk diagram: dua segmen dengan dasar yang berbeda menunjuk satu
alamat fisik yang sama.

\begin{figure}[H]
    \centering
    \begin{tikzpicture}[
        segk/.style={draw=obeBlue, line width=0.8pt, fill=obeBlueLight,
            rounded corners=2pt, minimum width=3.4cm, minimum height=1.0cm,
            align=center, font=\small, inner sep=4pt},
        offk/.style={draw=obeTeal, line width=0.8pt, fill=obeTealLight,
            rounded corners=2pt, minimum width=2.1cm, minimum height=0.85cm,
            align=center, font=\small\ttfamily, inner sep=4pt},
        fisk/.style={draw=obeAmber, line width=0.9pt, fill=obeAmberLight,
            rounded corners=2pt, minimum width=2.8cm, minimum height=1.3cm,
            align=center, font=\small, inner sep=4pt},
        pnrh/.style={-{Stealth[length=2.5mm]}, line width=0.8pt, obeBlue}
    ]
        \node[segk] (s1) at (0, 2.1)
            {Segmen \hex{0000}\\{\scriptsize dasar \hex{00000}}};
        \node[offk] (o1) at (3.7, 2.1) {\hex{0050}};

        \node[segk] (s2) at (0, 0.6)
            {Segmen \hex{0001}\\{\scriptsize dasar \hex{00010}}};
        \node[offk] (o2) at (3.7, 0.6) {\hex{0040}};

        \node[fisk] (f) at (7.6, 1.35)
            {Alamat fisik\\\hex{00050}};

        \draw[pnrh] (s1) -- (o1);
        \draw[pnrh] (s2) -- (o2);
        \draw[pnrh] (o1.east) -- ++(0.6,0) -| (f.north);
        \draw[pnrh] (o2.east) -- ++(0.6,0) -| (f.south);
    \end{tikzpicture}
    \caption{Dua pasangan \eng{segment:offset} yang berbeda menunjuk
    satu alamat fisik yang sama.}
    \label{fig:overlap-segmen}
\end{figure}

Berapa pasangan terbanyak yang mungkin menunjuk satu alamat?
Offset hanya mampu bergerak sejauh 64 KB, sedangkan setiap
kenaikan nomor segmen menggeser dasarnya 16 byte. Karena itu
jumlah pasangan untuk satu alamat tidak pernah melebihi
$65536 / 16 = 4096$ pasangan. Sebagai contoh, alamat \hex{10000}
dapat dinyatakan oleh 4096 pasangan, di antaranya
\code{0FFFh:0010h}, \code{0FFEh:0020h}, sampai
\code{1000h:0000h}. Secara keseluruhan terdapat
$65536 \times 65536 = \dec{4.294.967.296}$ alamat logis yang
berpadanan dengan hanya \dec{1.048.576} alamat fisik, sehingga
setiap alamat fisik rata-rata memiliki 4096 nama.

\begin{infobox}
\important{Banyak nilai segmen boleh sama, hasilnya tetap benar.}
Program tidak wajib memakai segmen yang berakhir nol atau segmen
yang ``rapi''. Selama pasangan \eng{segment:offset} menghasilkan
alamat fisik yang sama, isi memori yang diakses juga sama.
Karena itu dua program yang memakai pasangan berbeda tetap dapat
saling berbagi data, misalnya program yang menulis ke area video
lewat \code{B800h:0000h} dan program lain yang memakai
\code{B7FFh:0010h}.
\end{infobox}

\subsection{Akibat bagi Penulisan Program}

Konsep di atas menjelaskan mengapa 8086 memerlukan operator
\code{seg} di samping \code{OFFSET}. Operator \code{OFFSET}
menghasilkan bagian offset sebuah label, sedangkan operator
\code{seg} menghasilkan nomor segmennya. Kedua nilai itu bersama
membentuk sebuah \term{far pointer} berukuran 4 byte, yang di dalam
memori disimpan dengan urutan offset lebih dahulu, lalu segmen.
Pemakaian kedua operator tersebut ditunjukkan pada
Bagian~\ref{sec:contoh-program}, sedangkan aturan lengkap mengenai
segmen bawaan setiap mode pengalamatan dibahas pada
Bab~\ref{chap:mode-pengalamatan}.

Satu hal penting yang perlu diingat adalah bahwa pasangan
\eng{segment:offset} tidak boleh dipandang sebagai dua bagian yang
saling menggantikan. Perubahan kecil pada segmen menuntut perubahan
yang sepadan pada offset bila alamat fisik hendak dipertahankan.
Karena itu, ketika menelusuri variabel di dalam \eng{debugger},
sebaiknya catat kedua nilai tersebut sekaligus, bukan salah satu
saja.
